\section{Collection management and document lifecycle}
\label{sec:organization}

Collection membership determines where a document is listed and defines the
intended search scope when a collection identifier is supplied. The identifier
expresses requested scope rather than a hard access boundary: hierarchy
resolution and candidate expansion depend on the retrieval path. Parts of a
split document retain the collection assignment, and descriptions help the
agent select documents before opening their bodies. Table~\ref{tab:functions} pairs the management operations with their
platform or workflow checks.

\subsection{Classification, metadata, tags and duplicate review}
\label{sec:org-placement}
\label{sec:org-registries}
\label{sec:org-tags}
\label{sec:org-dedup}

The catalogue returns collections and sub-libraries with descriptions.
Each shelf shows document titles and descriptions within its selected scope.
The ingest agent compares the parsed text with collection descriptions
to choose a destination, then derives document descriptions and tags
from the body. Duplicate review
presents candidate pairs for inspection. Adjacent parts can describe
the same document while containing distinct source sections.
The placement and indexing steps connect managed documents to later query
access, whether through catalogue navigation or within a specified scope.

\begin{figure*}[!t]
\centering
\includegraphics[width=\textwidth]{fig_workflow_vector}
\caption{\looseness=-1 Illustrative HMRS workflow from document acquisition to evidence-based reporting. Documents are parsed when needed, organised into collections and indexed according to the entry point. A question can use catalogue navigation or pinned search; optional Laya scoring precedes evidence inspection. A completed review without support is distinct from an incomplete search or tool failure, which may occur at any stage.}
\label{fig:workflow}
\end{figure*}

HMRS supports library and sub-library creation, metadata edits, moves,
renames, deletion and forced reindexing. A move updates the document path
and registry entries, while reindexing updates the searchable representation.
\looseness=-1 Integrity checks compare document existence and path across the files on
disk, the \texttt{.tree-fs.json} registry and the library's
\texttt{.knowledge-base.yml}, and inspect index metadata.

Sub-libraries keep their descriptions and indices at their own tree nodes.
When the caller supplies a parent scope, the vector service resolves the
parent and its descendants into the collections searched. The lexical stage
of the two-stage service likewise resolves descendant identifiers. The
evaluation's flat arms (Section~\ref{sec:ladder-opt}) pass explicit collection identifiers rather than relying on a global
search call.

\subsection{Lifecycle operations, graph and operational knowledge}
\label{sec:org-lifecycle}
\label{sec:org-graph}
\label{sec:org-experience}

An early rename-and-move trial in a disposable collection exposed stale
chunks under the old document path until a forced reindex. The current rename
operation attempts to remove the old index entries and index the renamed
document; indexing failures are recorded separately from the metadata update.
A post-operation probe checks whether queries resolve to the new path.

The document graph links collections, documents and tags for navigation and
candidate expansion. Operational lessons are also stored as retrievable
documents. Experience tools track reviews, use and freshness and use
credibility tiers to order results. Section~\ref{sec:eval} reports
operational-experience cases alongside the management experiments.

\subsection{Organisation contract}
\label{sec:org-integrity}
\label{sec:org-summary}

The organisation contract records each document's current owning
collection, descriptive metadata and expected index representation. It does
not assert a one-to-one correspondence between a document and returned search
entries; selected probes are needed to detect duplicate, stale or unscoped
representations. For selected lifecycle operations, Section~\ref{sec:e2}
checks these relationships through file and metadata reconciliation,
description review and retrieval probes.
